\textbf{COLUNI 2018 -- Questão 09}\par\medskip Leia o excerto da carta – imaginária – escrita por Darcy Ribeiro endereçada à professora Diva Guimarães por ocasião de seu pronunciamento a respeito da situação social, histórica e política do negro no Brasil, na Flip de Paraty – RJ, em julho de 2017.

\par

De Darcy Ribeiro para a Diva da Flip

\par

“Precisamos, Diva, de muitas divas assim \textbf{como} você, que despertem uma crescente indignação contra esta herança maldita, porque é isso que nos dará forças para, amanhã, conter os possessos e criar aqui, neste país, uma sociedade solidária.” BRITO, Fernando. De Darcy Ribeiro para a Diva da FLIP. Disponível em: http://www.tijolaco.com.br/blog/de-darcy-ribeiro-para-diva-da flip/. Acesso em: 20 ago. 2017.

\par

O efeito de sentido construído, no excerto acima, pela conjunção \textbf{como} também está presente em:\begin{enumerate}
\item “Foi que eu vi as lagrimas deslizar dos olhos dos pobres. \textbf{Como} é pungente ver os dramas que ali se
       desenrola.” (Quarto de despejo, p. 42)
\item “Naquele momento, ela buscava na memória \textbf{como} o desenho da menina-flor tinha nascido em sua coleção.”
       (Zaíta esqueceu de guardar os brinquedos, p. 74)
\item “Estava, porém, chegando à conclusão de que trabalho \textbf{como} o dela não resolvia nada. Mas o que fazer?”
       (Zaíta esqueceu de guardar os brinquedos, p. 75)
\item “E o menino, Maria? \textbf{Como} vai o menino? cochichou o homem. Sabe que sinto falta de vocês? Tenho um
       buraco no peito, tamanha a saudade!” (Maria, p. 40)\end{enumerate}
